% Yunpeng Meng -- Curriculum Vitae
% Build: latexmk -pdf Yunpeng_Meng_CV.tex   (or run pdflatex twice)
% Red [ ... ] markers come from \todo{...}: replace each with the real fact, then delete the macro call.
\documentclass[11pt]{article}
\usepackage[letterpaper,top=0.7in,bottom=0.8in,left=0.95in,right=0.95in]{geometry}
\usepackage[T1]{fontenc}
\usepackage[utf8]{inputenc}
\usepackage{libertine}
\IfFileExists{newtxmath.sty}{\usepackage[libertine]{newtxmath}}{}
\usepackage[final]{microtype}
\usepackage{xcolor}
\usepackage{titlesec}
\usepackage{enumitem}
\usepackage{tabularx}
\usepackage{fancyhdr}
\usepackage{lastpage}
\definecolor{accent}{HTML}{27508F}   % same blue as the homepage
\definecolor{muted}{HTML}{5A6273}
\definecolor{todo}{HTML}{B3362F}
\usepackage[colorlinks=true,urlcolor=accent,linkcolor=accent,
  pdfauthor={Yunpeng Meng},pdftitle={Yunpeng Meng: Curriculum Vitae}]{hyperref}

\newcommand{\todo}[1]{\textcolor{todo}{\textsf{[#1]}}}

\setlength{\parindent}{0pt}
\setlength{\parskip}{0pt}
\pagestyle{fancy}
\fancyhf{}
\renewcommand{\headrulewidth}{0pt}
\fancyfoot[L]{\footnotesize\color{muted}Yunpeng Meng\enspace\textperiodcentered\enspace Curriculum Vitae}
\fancyfoot[R]{\footnotesize\color{muted}Updated October 2026\enspace\textperiodcentered\enspace Page \thepage\ of \pageref*{LastPage}}

\titleformat{\section}{\color{accent}\large\scshape}{}{0pt}{}[{\color{accent!35}\titlerule[0.6pt]}]
\titlespacing*{\section}{0pt}{1.1em}{0.55em}

% \cventry{left column: dates or label}{right column: details, lines separated by \newline}
\newcommand{\cventry}[2]{%
  \begin{tabularx}{\linewidth}{@{}>{\raggedright\arraybackslash}p{1.3in}@{}X@{}}
    \textcolor{muted}{#1} & #2
  \end{tabularx}\par\vspace{0.4em}}

\begin{document}

\begin{minipage}[b]{0.56\linewidth}
  {\fontsize{25}{29}\selectfont Yunpeng Meng}\\[4pt]
  {\large\color{muted}Curriculum Vitae}
\end{minipage}\hfill
\begin{minipage}[b]{0.42\linewidth}\raggedleft\small
  Department of Mathematics\\
  Texas A\&M University\\
  College Station, TX 77843, USA\\
  \href{mailto:meng.yun@tamu.edu}{meng.yun@tamu.edu}
\end{minipage}

\vspace{6pt}
{\color{accent}\rule{\linewidth}{0.9pt}}

\section{Education}
\cventry{Ph.D.}{\textbf{Texas A\&M University}, College Station, TX, USA\newline
  Mathematics, in progress\newline
  Advisor: Tian Yang}
\cventry{M.S.}{\textbf{Capital Normal University}, Beijing, China\newline
  Mathematics}
\cventry{B.S.}{\textbf{Tiangong University}, Tianjin, China\newline
  Mathematics}

\section{Research Interests}
Hyperbolic geometry and low-dimensional topology; the volume conjecture and quantum invariants of 3-manifolds;
geometric triangulations and combinatorial Ricci flow; asymptotics of quantum $6j$-symbols;
resurgence and Borel--Laplace resummation.

\section{Papers and Preprints}
\begin{enumerate}[leftmargin=2.4em,labelsep=0.7em,itemsep=0.55em,topsep=0pt,parsep=0pt]
  \item[{[4]}] \textit{On the volume conjectures for Dehn-filled 3-manifolds along fundamental shadow link
    complements with big volume}, with Baojun Wu, Tian Yang, and Yuxuan Yang.\newline
    Preprint, 2026.
  \item[{[3]}] \textit{Asymptotics of complex $b$-$6j$ symbols}, with Tian Yang.\newline
    \href{https://arxiv.org/abs/2606.03653}{arXiv:2606.03653}, 2026.
  \item[{[2]}] \textit{Hyperbolization and geometric decomposition of a class of 3-manifolds},
    with Ke Feng and Huabin Ge.\newline
    \href{https://arxiv.org/abs/2503.07421}{arXiv:2503.07421}, 2025.
  \item[{[1]}] \textit{On the volume conjecture for hyperbolic Dehn-filled 3-manifolds along the twist knots},
    with Huabin Ge, Chuwen Wang, and Yuxuan Yang.\newline
    \href{https://arxiv.org/abs/2410.20489}{arXiv:2410.20489}, 2024.
\end{enumerate}

\section{Teaching}
\cventry{Fall 2026}{Grader, MATH 415H \textit{Modern Algebra I} (Honors), Texas A\&M University\newline
  Instructor: Frank Sottile}

\section{References}
\cventry{}{\textbf{Tian Yang}, Ph.D.\ advisor\newline
  Associate Professor, Department of Mathematics\newline
  Texas A\&M University, College Station, TX 77843, USA}

\end{document}
